The given sky map (in Mercator projection) shows the constellation of Orion as
seen from a certain location X (longitude $\lambda_X = 70^\circ$\,E) on
21 March 2025 at 22:00 local time. The point ``W'' marks the cardinal West
point. The altitude and azimuth values are marked on the grid.

% FIGURE NEEDED: Mercator-projection sky map of Orion with altitude/azimuth
% grid, points A-F and cardinal point W marked (2025_SKYMAP paper)

\begin{description}
\item[(OM03.1)] What is the approximate latitude ($\phi_X$) of the
  location X?
\item[(OM03.2)] The grid provided on the Summary Answersheet has the same
  projection (Mercator), and the angular scale in both azimuth and altitude
  are identical to that of the grid provided in the question. On this grid
  draw to scale the constellation of Orion as it will be at another location Y
  (with latitude $\phi_Y = 40^\circ$\,S and longitude
  $\lambda_Y = 50^\circ$\,W) on 21 January 2026 at 18:00 local time. An
  approximate outline of the constellation is enough, with the points A--F
  being marked clearly. Identify the cardinal point ``P'' shown on the grid
  (tick ($\checkmark$) the appropriate box in the Summary Answersheet). You
  may make suitable approximations to arrive at your answer.

  You may use the following relations between the Hour Angle ($H$),
  Declination ($\delta$), Altitude ($a$), Azimuth ($A$) and Latitude
  ($\phi$):
  \[ \cos H = \frac{\sin a - \sin\delta \sin\phi}{\cos\delta \cos\phi} \]
  \[ \sin\delta = \sin\phi \sin a + \cos\phi \cos a \cos A \]
\end{description}
